that are isomorphisms at the points @@M0@@. Now consider the complex:
@@M1@@
The class of @@M2@@ is @@M3@@; moreover, @@M4@@ is acyclic at the points @@M5@@, so @@M6@@ contains no irreducible component of @@M7@@ and consequently has codimension @@M8@@ in @@M9@@; this shows that @@M10@@ belongs to @@M11@@.

Now let @@M12@@ be an element of @@M13@@; we shall show that @@M14@@. We may suppose that @@M15@@ can be written:
@@M16@@
the @@M17@@ being the classes of locally free sheaves of finite type and, for every @@M18@@, @@M19@@ and @@M20@@ having the same rank. Moreover, decomposing @@M21@@ into a sum of finitely many open sets and arguing in each open set, we may suppose the @@M22@@ to have constant rank. Let @@M23@@ be the fibre product over @@M24@@ of the flag schemes of the @@M25@@. In @@M26@@, the images of the @@M27@@ are sums of classes of invertible sheaves. Let @@M28@@ be the classes of two invertible @@M29@@-Modules; we have:
@@M30@@
Thus @@M31@@. Now @@M32@@ is noetherian and has an ample sheaf, so, by what precedes and 6.6, @@M33@@. Using the additivity of @@M34@@ and 6.6, we deduce that @@M35@@, and @@M36@@, where @@M37@@ is the canonical morphism @@M38@@.

We shall deduce that @@M39@@. Let @@M40@@ be the relative dimension of @@M41@@ over @@M42@@. There is an element @@M43@@ of @@M44@@, belonging to a basis of @@M45@@ over @@M46@@, such that:
